% =============================================================================
%  1-Objetivos.tex — Objetivos de la sesión
% =============================================================================

\section{Objetivos de la sesión}

Al finalizar la sesión, se espera que quienes participen:

\begin{enumerate}
  \item Se aproximen, mediante una experiencia inmersiva breve y controlada dentro del aula, a algunas de las barreras sensoriales y cognitivas que enfrentan las personas autistas, para comprender el autismo desde la vivencia y no solo desde la definición.
  \item Identifiquen y distingan los conceptos de enmascaramiento (\textit{masking}), hiperfoco, burnout, meltdown y shutdown, a partir de las descripciones de las propias personas autistas.
  \item Reconozcan los estímulos del entorno, como el ruido, la iluminación, la carga visual y las reglas sociales implícitas, que pueden limitar el aprendizaje y la participación, tanto en el aula como en el espacio urbano.
  \item Reflexionen críticamente sobre su papel como urbanistas y docentes en el diseño de espacios y dinámicas que reconozcan la neurodiversidad.
\end{enumerate}
